\documentclass[11pt]{article}
\usepackage[margin=1in]{geometry}
\usepackage{xcolor}
\usepackage{tcolorbox}
\usepackage{hyperref}
\usepackage{enumitem}
\setlist{noitemsep,topsep=0pt}
\usepackage{booktabs}
\usepackage{array}
\usepackage{amsmath}

% Define OSU orange and AI blue colors
\definecolor{osaorange}{RGB}{220,115,38}
\definecolor{aiblue}{RGB}{30,144,255}

% Configure hyperref colors
\hypersetup{
    colorlinks=true,
    linkcolor=blue,
    urlcolor=blue,
    citecolor=blue
}

\title{}
\author{}
\date{}

\begin{document}

% Orange header box with black outline
\begin{tcolorbox}[colback=osaorange,colframe=black,boxrule=2pt,arc=0pt,left=6pt,right=6pt,top=10pt,bottom=10pt]
\centering
\textcolor{white}{\textbf{\Large In-Class Worksheet - ANSWER KEY}}\\
\textcolor{white}{\textbf{\large Probability and Conditional Probability}}
\end{tcolorbox}

\vspace{8pt}

\noindent\textbf{Course:} H524 Introduction to Biostatistics\\
\textbf{Topic:} Probability, Conditional Probability, and Diagnostic Testing

\vspace{8pt}

\section{Problem 1: Mammogram Screening (Based on Gigerenzer's Study)}

A 45-year-old woman with no symptoms or family history of breast cancer has a positive mammogram. What is the probability that she actually has breast cancer?

\vspace{6pt}

\textbf{Given information:}
\begin{itemize}
\item The prevalence of breast cancer in this age group is 1.2\%
\item If a woman has breast cancer, the probability is 85\% that she will have a positive mammogram (sensitivity)
\item If a woman does not have breast cancer, the probability is 9\% that she will still have a positive mammogram (false positive rate)
\end{itemize}

\vspace{12pt}

\textbf{Part A:} Set up a 2x2 table assuming 10,000 women are screened. Fill in all cells:

\vspace{12pt}

\begin{center}
\begin{tabular}{|l|c|c|c|}
\hline
& \textbf{Cancer +} & \textbf{Cancer -} & \textbf{Total} \\
\hline
\textbf{Test +} & 102 & 889 & 991 \\[2ex]
\hline
\textbf{Test -} & 18 & 8,991 & 9,009 \\[2ex]
\hline
\textbf{Total} & 120 & 9,880 & 10,000 \\[2ex]
\hline
\end{tabular}
\end{center}

\vspace{12pt}

\textbf{Work space for calculations:}

\begin{itemize}
\item Total with cancer: $10,000 \times 0.012 = 120$
\item Total without cancer: $10,000 - 120 = 9,880$
\item True positives: $120 \times 0.85 = 102$
\item False negatives: $120 - 102 = 18$
\item False positives: $9,880 \times 0.09 = 889$
\item True negatives: $9,880 - 889 = 8,991$
\end{itemize}

\vspace{12pt}

\textbf{Part B:} Calculate the Positive Predictive Value (PPV): $\Pr(\text{Cancer} \mid \text{Test+})$

$$\text{PPV} = \frac{102}{991} = 0.103 = \textbf{10.3\%}$$

\vspace{6pt}

\textbf{Part C:} Interpret your result. Why is the PPV so much lower than the sensitivity?

\textbf{Answer:} Even though the test has 85\% sensitivity, only 10.3\% of positive tests indicate actual cancer. This is because the prevalence is low (1.2\%), so false positives (889) vastly outnumber true positives (102). In a low-prevalence population, even a small false positive rate (9\%) generates many false positives.

\newpage

\section{Problem 2: Strep Throat Testing}

A rapid strep test has 85\% sensitivity and 95\% specificity. Among patients presenting with sore throat, 15\% actually have strep throat.

\vspace{12pt}

\textbf{Part A:} Create a 2x2 table for 1,000 patients with sore throat:

\vspace{12pt}

\begin{center}
\begin{tabular}{|l|c|c|c|}
\hline
& \textbf{Strep +} & \textbf{Strep -} & \textbf{Total} \\
\hline
\textbf{Test +} & 128 & 43 & 171 \\[2ex]
\hline
\textbf{Test -} & 22 & 807 & 829 \\[2ex]
\hline
\textbf{Total} & 150 & 850 & 1,000 \\[2ex]
\hline
\end{tabular}
\end{center}

\textbf{Calculations:}
\begin{itemize}
\item Strep positive: $1,000 \times 0.15 = 150$
\item Strep negative: $1,000 - 150 = 850$
\item True positives: $150 \times 0.85 = 128$ (rounded)
\item False negatives: $150 - 128 = 22$
\item True negatives: $850 \times 0.95 = 808$ (rounded to 807 for total)
\item False positives: $850 - 807 = 43$
\end{itemize}

\vspace{12pt}

\textbf{Part B:} Calculate the Negative Predictive Value (NPV): $\Pr(\text{No Strep} \mid \text{Test-})$

$$\text{NPV} = \frac{807}{829} = 0.973 = \textbf{97.3\%}$$

\vspace{12pt}

\textbf{Part C:} What is the probability that a patient with a negative test actually has strep (false negative)?

$$\Pr(\text{Strep} \mid \text{Test-}) = \frac{22}{829} = 0.027 = \textbf{2.7\%}$$

\section{Problem 3: Smoking and Heart Disease}

In a study of 1,000 people, there were 200 smokers and 800 non-smokers. During the 10-year follow-up, 30 smokers and 20 non-smokers developed heart disease.

\vspace{12pt}

\textbf{Part A:} Create a 2x2 table:

\vspace{12pt}

\begin{center}
\begin{tabular}{|l|c|c|c|}
\hline
& \textbf{Heart Disease} & \textbf{No Disease} & \textbf{Total} \\
\hline
\textbf{Smokers} & 30 & 170 & 200 \\[2ex]
\hline
\textbf{Non-smokers} & 20 & 780 & 800 \\[2ex]
\hline
\textbf{Total} & 50 & 950 & 1,000 \\[2ex]
\hline
\end{tabular}
\end{center}

\vspace{12pt}

\textbf{Part B:} Calculate the risk of heart disease in smokers:

$$\text{Risk}_{\text{smokers}} = \frac{30}{200} = 0.15 = \textbf{15\%}$$

\vspace{6pt}

\textbf{Part C:} Calculate the risk of heart disease in non-smokers:

$$\text{Risk}_{\text{non-smokers}} = \frac{20}{800} = 0.025 = \textbf{2.5\%}$$

\vspace{6pt}

\textbf{Part D:} Calculate the Relative Risk (RR):

$$\text{RR} = \frac{\text{Risk}_{\text{smokers}}}{\text{Risk}_{\text{non-smokers}}} = \frac{0.15}{0.025} = \textbf{6.0}$$

\vspace{6pt}

\textbf{Part E:} Interpret the relative risk in plain language:

\textbf{Answer:} Smokers have 6 times the risk of developing heart disease compared to non-smokers over the 10-year period. Smoking is strongly associated with increased heart disease risk.

\section{Problem 4: Binomial Probability}

A new treatment has a 60\% success rate.

\vspace{12pt}

\textbf{Part A:} If 10 patients receive the treatment, what is the probability that at least 8 of them respond successfully? Use R or a binomial probability table.

\textbf{Using R:}
\begin{verbatim}
pbinom(7, size = 10, prob = 0.6, lower.tail = FALSE)
\end{verbatim}

$$\Pr(X \geq 8) = \Pr(X = 8) + \Pr(X = 9) + \Pr(X = 10) = \textbf{0.167 or 16.7\%}$$

\vspace{12pt}

\textbf{Part B:} What is the expected number of successful responses out of 10 patients?

$$E(X) = n \times p = 10 \times 0.6 = \textbf{6 patients}$$

\section{Problem 5: Normal Distribution}

Cholesterol levels in adults are normally distributed with mean 200 mg/dL and standard deviation 40 mg/dL.

\vspace{12pt}

\textbf{Part A:} What proportion of adults have cholesterol levels between 180 and 220 mg/dL?

\textbf{Solution:}
\begin{itemize}
\item $Z_1 = \frac{180 - 200}{40} = -0.5$
\item $Z_2 = \frac{220 - 200}{40} = 0.5$
\item $\Pr(-0.5 < Z < 0.5) = \Pr(Z < 0.5) - \Pr(Z < -0.5) = 0.6915 - 0.3085 = \textbf{0.383 or 38.3\%}$
\end{itemize}

\textbf{Using R:}
\begin{verbatim}
pnorm(220, mean = 200, sd = 40) - pnorm(180, mean = 200, sd = 40)
\end{verbatim}

\vspace{12pt}

\textbf{Part B:} What cholesterol level corresponds to the 90th percentile?

\textbf{Solution:}
\begin{itemize}
\item The 90th percentile has $Z = 1.28$ (from Z-table)
\item $X = \mu + Z\sigma = 200 + 1.28(40) = 200 + 51.2 = \textbf{251.2 mg/dL}$
\end{itemize}

\textbf{Using R:}
\begin{verbatim}
qnorm(0.90, mean = 200, sd = 40)
\end{verbatim}

\end{document}
